%==============================================================
%  lecons/algebre/corps_fractions.tex — Corps des fractions
%==============================================================

\lecontitle{Corps des fractions}{Algèbre — Anneaux intègres}

%=============================================================
%  § 1 — PRÉREQUIS : ANNEAU INTÈGRE
%=============================================================
\secnum{navyblue}{1}{Anneaux intègres — rappels et motivation}

\begin{tcolorbox}[thmbox]
  \textbf{Définition.}\enspace%
  \index{anneau intègre}
  Un anneau commutatif $A$ (avec $1 \neq 0$) est dit \emph{intègre} s'il
  n'a pas de \emph{diviseur de zéro} :
  \[
    ab = 0 \;\Rightarrow\; a = 0 \text{ ou } b = 0.
  \]
  \textbf{Exemples :} $\mathbb{Z}$, $\mathbb{K}[X]$ pour $\mathbb{K}$ corps,
  tout corps, $\mathbb{Z}[i]$ (entiers de Gauss).\\
  \textbf{Contre-exemples :} $\mathbb{Z}/6\mathbb{Z}$ ($2\cdot 3=0$),
  $\mathcal{M}_2(\mathbb{K})$.
\end{tcolorbox}

\rubriq{Motivation}
Dans $\mathbb{Z}$, on peut additionner et multiplier, mais pas toujours diviser :
$3/2 \notin \mathbb{Z}$. Le corps des fractions est la plus petite extension
dans laquelle tout élément non nul de $A$ est inversible — on « force » la division.

\decsep{}

%=============================================================
%  § 2 — CONSTRUCTION
%=============================================================
\secnum{navyblue}{2}{Construction du corps des fractions}

\begin{tcolorbox}[thmbox]
  \textbf{Théorème-définition.}\enspace%
  \index{corps des fractions}
  Soit $A$ un anneau intègre. Sur $S = A \times (A\setminus\{0\})$, on définit
  \[
    (a,b) \sim (c,d) \;\iff\; ad = bc.
  \]
  \textit{C'est une relation d'équivalence. L'ensemble quotient
  $\mathrm{Frac}(A) = S/{\sim}$, muni des lois}
  \[
    \frac{a}{b} + \frac{c}{d} = \frac{ad+bc}{bd},
    \qquad
    \frac{a}{b} \cdot \frac{c}{d} = \frac{ac}{bd},
  \]
  \textit{est un corps, appelé \emph{corps des fractions} de $A$.
  L'application $\iota : A \to \mathrm{Frac}(A),\; a \mapsto \tfrac{a}{1}$,
  est un homomorphisme d'anneaux injectif.}
\end{tcolorbox}

\rubriq{Preuve — $\sim$ est bien une relation d'équivalence}

\begin{mdframed}[style=proofbar]
  \textit{Réflexivité.} $ab = ba$ (commutativité). $\checkmark$

  \textit{Symétrie.} $ad=bc \Rightarrow cb=da$. $\checkmark$

  \textit{Transitivité.} Soit $(a,b)\sim(c,d)$ et $(c,d)\sim(e,f)$
  (où $b,d,f$ sont tous non nuls par construction).
  Supposons donc $ad=bc$ et $cf=de$.
  On veut $af=be$. En multipliant la première égalité par $f$,
  $adf = bcf = b(cf) = b(de) = bde$,
  donc $d(af-be)=0$. Comme $d\neq 0$ et $A$ est \emph{intègre}, on peut
  simplifier par $d$ : $af=be$. $\blacksquare$
\end{mdframed}

\rubriq{Preuve — $\mathrm{Frac}(A)$ est un corps}

\begin{mdframed}[style=proofbar]
  \textit{Les lois ont un sens.} Comme $b,d\neq 0$ et que $A$ est intègre,
  $bd\neq 0$ : les seconds membres $\tfrac{ad+bc}{bd}$ et $\tfrac{ac}{bd}$ sont
  bien des éléments de $S$ (c'est encore l'intégrité qui assure que le
  dénominateur reste non nul).

  \smallskip
  \textit{Bien-fondé des lois.} Si $(a,b)\sim(a',b')$ et $(c,d)\sim(c',d')$
  (i.e.\ $ab'=a'b$ et $cd'=c'd$), montrons $(ad+bc,\,bd)\sim(a'd'+b'c',\,b'd')$ :
  \begin{align*}
    (ad+bc)(b'd') &= adb'd' + bcb'd' = (ab')dd' + (cd')bb' \\
                  &= (a'b)dd'+(c'd)bb' = (a'd'+b'c')(bd).
  \end{align*}
  Pour le produit, de même
  $(ac)(b'd') = (ab')(cd') = (a'b)(c'd) = (a'c')(bd)$, donc
  $(ac,\,bd)\sim(a'c',\,b'd')$.

  \textit{Anneau commutatif avec $1=\tfrac{1}{1}$ et $0=\tfrac{0}{1}$.}
  Associativité, commutativité et distributivité se ramènent aux propriétés
  correspondantes de $A$ après réduction au même dénominateur ; par exemple
  $\tfrac{a}{b}+\tfrac{0}{1}=\tfrac{a\cdot 1+b\cdot 0}{b\cdot 1}=\tfrac{a}{b}$ et
  $\tfrac{a}{b}\cdot\tfrac{1}{1}=\tfrac{a}{b}$, ce qui confirme les éléments neutres.

  \textit{Inverse de $\tfrac{a}{b}$ pour $a\neq 0$.}
  $\dfrac{a}{b}\cdot\dfrac{b}{a} = \dfrac{ab}{ba} = \dfrac{1}{1}$.
  Ainsi tout élément non nul est inversible : $\mathrm{Frac}(A)$ est un corps.

  \textit{Injectivité de $\iota$.}
  $\iota(a)=\iota(a') \Rightarrow \tfrac{a}{1}=\tfrac{a'}{1}
  \Rightarrow (a,1)\sim(a',1) \Rightarrow a\cdot 1 = a'\cdot 1 \Rightarrow a=a'$. $\blacksquare$
\end{mdframed}

\decsep{}

%=============================================================
%  § 3 — PROPRIÉTÉ UNIVERSELLE
%=============================================================
\secnum{navyblue}{3}{Propriété universelle}

\begin{tcolorbox}[thmbox]
  \textbf{Théorème (propriété universelle).}\enspace%
  \index{corps des fractions!propriété universelle}
  \textit{Soit $\varphi : A \to K$ un homomorphisme d'anneaux injectif vers
  un corps $K$. Alors il existe un unique homomorphisme de corps
  $\tilde{\varphi} : \mathrm{Frac}(A) \to K$ tel que $\tilde{\varphi}\circ\iota = \varphi$,
  à savoir}
  \[
    \tilde{\varphi}\!\left(\frac{a}{b}\right) = \frac{\varphi(a)}{\varphi(b)}.
  \]
  \textit{En particulier, $\mathrm{Frac}(A)$ est, à isomorphisme près, le plus
  petit corps contenant $A$ (via le plongement $\iota$).}
\end{tcolorbox}

\begin{mdframed}[style=proofbar]
  \textit{Bien-fondé.} Si $(a,b)\sim(c,d)$, alors $ad=bc$, donc
  $\varphi(a)\varphi(d)=\varphi(b)\varphi(c)$.
  Comme $\varphi$ est injectif et $b\neq 0$, $\varphi(b)\neq 0$, d'où
  $\varphi(a)/\varphi(b) = \varphi(c)/\varphi(d)$ dans $K$.

  \textit{Homomorphisme.} Vérification directe sur la somme et le produit.

  \textit{Unicité.} Comme $\tfrac{b}{1}\cdot\tfrac{1}{b}=1$ dans
  $\mathrm{Frac}(A)$, on a nécessairement
  $\tilde\varphi(1/b)={\tilde\varphi(b/1)}^{-1}={\varphi(b)}^{-1}$ ;
  $\tilde\varphi$ est donc forcée par $\tilde\varphi(a/b)
  = \tilde\varphi(a/1)\cdot\tilde\varphi(1/b)
  = \varphi(a)\cdot{\varphi(b)}^{-1}$. $\blacksquare$
\end{mdframed}

\decsep{}

%=============================================================
%  § 4 — EXEMPLES
%=============================================================
\secnum{navyblue}{4}{Exemples fondamentaux}

\begin{mdframed}[style=exbar]
  \textbf{1. $\mathrm{Frac}(\mathbb{Z}) = \mathbb{Q}$.}\enspace%
  \index{nombres rationnels!construction}
  La construction ci-dessus avec $A=\mathbb{Z}$ redonne exactement les rationnels,
  avec $\tfrac{a}{b}$ au sens habituel — ce cas particulier est détaillé dans
  la leçon consacrée à la construction de $\mathbb{Q}$.

  \smallskip\noindent
  \textbf{2. Corps des fractions rationnelles $\mathbb{K}(X)$.}\enspace%
  \index{fractions rationnelles}
  $\mathrm{Frac}(\mathbb{K}[X]) = \mathbb{K}(X)$, les \emph{fractions rationnelles}
  $P/Q$ avec $P,Q\in\mathbb{K}[X]$, $Q\neq 0$.
  Plus généralement, $\mathrm{Frac}(\mathbb{K}[X_1,\ldots,X_n]) = \mathbb{K}(X_1,\ldots,X_n)$.

  \smallskip\noindent
  \textbf{3. Entiers de Gauss.}\enspace%
  \index{entiers de Gauss}
  $\mathrm{Frac}(\mathbb{Z}[i]) = \mathbb{Q}(i) = \{a+bi \mid a,b\in\mathbb{Q}\}$,
  le corps des \emph{rationnels de Gauss}.

  \smallskip\noindent
  \textbf{4. Corps de fonctions d'une courbe.}\enspace%
  Si $A = \mathbb{K}[X,Y]/(f)$ est l'anneau de coordonnées d'une courbe algébrique
  irréductible $f=0$, alors $\mathrm{Frac}(A)$ est le corps de fonctions de la courbe,
  objet central de la géométrie algébrique.
\end{mdframed}

\decsep{}

%=============================================================
%  § 5 — PIÈGES ET EXERCICES
%=============================================================
\secnum{exgreen}{5}{Pièges et exercices}

\rubriq{Pièges}

\begin{mdframed}[style=cexbar]
  \textbf{L'intégrité est indispensable pour la transitivité.}\enspace%
  Dans $\mathbb{Z}/6\mathbb{Z}$, où la relation s'écrit
  $(a,b)\sim(c,d)\iff ad=bc$ dans $\mathbb{Z}/6\mathbb{Z}$
  — soit, en raisonnant sur des représentants entiers,
  $ad\equiv bc\pmod 6$ —, on a
  $(2,2)\sim(3,3)$ car $2\cdot 3=6\equiv 0$ et $2\cdot 3=6\equiv 0\pmod 6$, et
  $(3,3)\sim(2,4)$ car $3\cdot 4=12\equiv 0$ et $3\cdot 2=6\equiv 0\pmod 6$ ;
  pourtant $(2,2)\not\sim(2,4)$ car $2\cdot 4=8\equiv 2$ tandis que
  $2\cdot 2=4\not\equiv 2\pmod 6$.
  La transitivité tombe en défaut : la construction échoue sur un anneau non intègre.

  \smallskip\noindent
  \textbf{$\iota$ n'est pas surjective en général.}\enspace%
  $\iota : \mathbb{Z} \to \mathbb{Q}$ n'est pas surjective : $1/2 \notin \mathrm{Im}(\iota)$.
  $\mathrm{Frac}(A)$ est strictement plus grand que $A$ dès que $A$ n'est pas
  déjà un corps.

  \smallskip\noindent
  \textbf{$\mathrm{Frac}(A) = A$ si et seulement si $A$ est un corps.}\enspace%
  Dans ce cas, $\iota$ est un isomorphisme.
\end{mdframed}

\vspace{6pt}
\exolabel

\begin{mdframed}[style=exobar]
  \begin{enumerate}
    \item Montrer que si $A$ est intègre et $\mathrm{Frac}(A) = A$,
          alors $A$ est un corps.

    \item Déterminer $\mathrm{Frac}(\mathbb{Z}[\sqrt{2}])$.
          Montrer que c'est $\mathbb{Q}(\sqrt{2}) = \{a+b\sqrt{2}\mid a,b\in\mathbb{Q}\}$.

    \item Soit $p$ premier. On note $\mathbb{Z}_{(p)}$ l'ensemble des rationnels
          qui \emph{peuvent s'écrire} $\tfrac{a}{b}$ avec $p\nmid b$
          (vérifier que cette condition équivaut à \og $p$ ne divise pas le
          dénominateur de la forme irréductible \fg{}, donc ne dépend pas du
          représentant choisi). Montrer que $\mathbb{Z}_{(p)}$ est un sous-anneau
          de $\mathbb{Q}$ (l'anneau local en $p$), et que son
          unique idéal maximal est $p\mathbb{Z}_{(p)}$.

    \item Décomposer en éléments simples dans $\mathbb{R}(X)$ la fraction
          $\dfrac{X^2+1}{(X-1)(X^2+X+1)}$.

    \item Soit $A \subset B$ des anneaux intègres.
          Montrer que $\mathrm{Frac}(A) \subset \mathrm{Frac}(B)$.
          Si $\mathrm{Frac}(A) = \mathrm{Frac}(B)$, a-t-on $A = B$ ?
          (Donner un contre-exemple.)
  \end{enumerate}
\end{mdframed}

\leconfooter{R.\ Dedekind (1871) — H.\ Weber (1882) — fondements de l'algèbre commutative}
